\documentclass[11pt,letterpaper]{article}
\usepackage[T1]{fontenc}
\usepackage[utf8]{inputenc}
\usepackage[margin=0.88in,top=0.84in,bottom=0.83in,headheight=13pt,headsep=20pt,footskip=28pt]{geometry}
\usepackage{newpxtext}
\usepackage{amsmath,amsthm}
\usepackage{newpxmath}
\usepackage{microtype}
\usepackage{xcolor}
\definecolor{Forest}{HTML}{30392C}
\definecolor{Olive}{HTML}{687144}
\definecolor{Muted}{HTML}{65685F}
\definecolor{Sage}{HTML}{CBD0BC}
\definecolor{Pale}{HTML}{F2F4EB}
\usepackage{mathtools}
\usepackage{booktabs,tabularx,longtable,array}
\usepackage{enumitem}
\usepackage{titlesec}
\usepackage{fancyhdr}
\usepackage[most]{tcolorbox}
\usepackage{needspace}
\PassOptionsToPackage{obeyspaces}{url}
\usepackage{xurl}
\usepackage[colorlinks=true,linkcolor=Olive,citecolor=Olive,urlcolor=Olive,bookmarksnumbered=true,pdfencoding=auto,psdextra]{hyperref}
\hypersetup{pdftitle={Normal measures from critical-sequence tail classes},pdfauthor={Research continuation prepared for Vladimir Reshetnikov},pdfsubject={Ultraexacting cardinals, relative HOD, internal normal measures, Prikry genericity, and conditional consistency}}
\urlstyle{same}
\setlength{\parindent}{0pt}
\setlength{\parskip}{5.1pt plus 1pt minus .4pt}
\linespread{1.035}
\setlength{\emergencystretch}{2em}
\setcounter{tocdepth}{1}
\setcounter{secnumdepth}{2}
\titleformat{\section}{\Large\bfseries\color{Forest}}{\thesection}{.6em}{}
\titleformat{\subsection}{\large\bfseries\color{Forest}}{\thesubsection}{.55em}{}
\titleformat{\subsubsection}{\normalsize\bfseries\color{Forest}}{}{0pt}{}
\titlespacing*{\section}{0pt}{18pt plus 3pt minus 2pt}{8pt}
\titlespacing*{\subsection}{0pt}{12pt plus 2pt minus 1pt}{5pt}
\titlespacing*{\subsubsection}{0pt}{9pt}{3pt}
\setlist[itemize]{leftmargin=1.4em,itemsep=2pt,topsep=3pt,parsep=0pt}
\setlist[enumerate]{leftmargin=1.7em,itemsep=3pt,topsep=4pt,parsep=0pt}
\pagestyle{fancy}
\fancyhf{}
\fancyhead[L]{\sffamily\fontsize{9}{11}\selectfont\color{Muted}CRITICAL-SEQUENCE TAIL CLASSES}
\fancyhead[R]{\sffamily\fontsize{9}{11}\selectfont\color{Muted}CONTINUATION\quad 18 SEP 2026}
\fancyfoot[L]{\small\color{Muted}Normal measures and canonical Prikry layers}
\fancyfoot[R]{\small\color{Muted}\thepage}
\renewcommand{\headrulewidth}{.35pt}
\renewcommand{\footrulewidth}{0pt}
\fancypagestyle{plain}{\fancyhf{}\fancyfoot[L]{\small\color{Muted}Normal measures and canonical Prikry layers}\fancyfoot[R]{\small\color{Muted}\thepage}\renewcommand{\headrulewidth}{0pt}}
\tcbset{colback=Pale,colframe=Sage,colbacktitle=Sage,coltitle=Forest,fonttitle=\bfseries,boxrule=.5pt,arc=3pt,left=9pt,right=9pt,top=6pt,bottom=6pt,toptitle=3pt,bottomtitle=3pt,before skip=8pt,after skip=8pt,breakable}
\newtcolorbox{assessment}[1]{title={#1}}
\theoremstyle{definition}
\newtheorem{definition}{Definition}[section]
\newtheorem{theorem}[definition]{Theorem}
\newtheorem{proposition}[definition]{Proposition}
\newtheorem{lemma}[definition]{Lemma}
\newtheorem{corollary}[definition]{Corollary}
\newtheorem{remark}[definition]{Remark}
\newtheorem{question}[definition]{Question}
\newtheorem{inputthm}{Input}
\renewcommand{\theinputthm}{\Alph{inputthm}}
\numberwithin{equation}{section}
\newcommand{\ZFC}{\mathsf{ZFC}}
\newcommand{\ZF}{\mathsf{ZF}}
\newcommand{\AC}{\mathsf{AC}}
\newcommand{\GA}{\mathsf{GA}}
\newcommand{\HOD}{\mathrm{HOD}}
\newcommand{\OD}{\mathrm{OD}}
\newcommand{\CD}{\mathrm{CD}}
\newcommand{\HCD}{\mathrm{HCD}}
\newcommand{\UF}{\mathrm{UF}}
\newcommand{\Ex}{\operatorname{Ex}}
\newcommand{\UEx}{\operatorname{UEx}}
\newcommand{\Ext}{\operatorname{Ext}}
\newcommand{\SC}{\operatorname{SC}}
\newcommand{\CEx}{\operatorname{CEx}}
\newcommand{\REx}{\operatorname{REx}}
\newcommand{\Con}{\operatorname{Con}}
\newcommand{\crit}{\operatorname{crit}}
\newcommand{\cf}{\operatorname{cf}}
\newcommand{\ot}{\operatorname{ot}}
\newcommand{\ran}{\operatorname{ran}}
\newcommand{\rank}{\operatorname{rank}}
\newcommand{\lcm}{\operatorname{lcm}}
\newcommand{\Ord}{\mathrm{Ord}}
\newcommand{\Pow}{\mathcal P}
\newcommand{\id}{\mathrm{id}}
\newcommand{\restr}{\mathbin{\upharpoonright}}
\newcommand{\Izero}{\mathrm{I0}}
\newcommand{\Itwo}{\mathrm{I2}}
\newcommand{\Ithree}{\mathrm{I3}}
\newcommand{\Iwf}{\mathrm{I3}_{\mathrm{wf}}(0)}
\newcommand{\Sel}{\operatorname{Sel}}
\newcommand{\FF}{\mathcal F}
\newcommand{\PP}{\mathbb P}
\newcommand{\Clam}{\mathcal C_\lambda}
\newcommand{\Rig}{\operatorname{Rig}}
\newcommand{\Spec}{\operatorname{Spec}}
\newcommand{\ch}{\operatorname{ind}}
\newcommand{\CF}{\mathsf{CF}}
\newcommand{\src}[1]{\unskip\ {\normalfont\small\sffamily\color{Olive}[#1]}\ }
\newcommand{\R}[1]{\textsf{R#1}}
\newcolumntype{Y}{>{\raggedright\arraybackslash}X}
\newcolumntype{P}[1]{>{\raggedright\arraybackslash}p{#1}}
\newcolumntype{C}{>{\centering\arraybackslash}p{.034\textwidth}}
\newcommand{\yes}{$\bullet$}
\newcommand{\prt}{$\circ$}
\newcommand{\arxiv}[1]{\href{https://arxiv.org/abs/#1}{arXiv:#1}}
\newcommand{\doi}[1]{\href{https://doi.org/#1}{doi:\nolinkurl{#1}}}


\newcommand{\tc}{\operatorname{tc}}
\newcommand{\Tail}{\operatorname{Tail}}
\newcommand{\Split}{\operatorname{Split}}
\newcommand{\Mqp}{M_{q,p}}
\newcommand{\Uqp}{U_{q,p}}
\newcommand{\Kqp}{K_{q,p}}
\newcommand{\code}[1]{\ulcorner #1\urcorner}
\newcommand{\Refl}{\operatorname{Refl}_{\omega}}
\newcommand{\NS}{\mathsf{DSplit}}
\newcommand{\NMPr}{\mathsf{NMPr}}
\newcommand{\suff}{\mathbin{\smallfrown}}

\begin{document}
\thispagestyle{empty}
{\sffamily\bfseries\small\color{Olive}SET THEORY\quad /\quad FOURTH RESEARCH CONTINUATION}

\vspace{13pt}
{\fontsize{28}{31}\selectfont\bfseries\color{Forest}Normal measures from\\critical-sequence tail classes\par}
\vspace{10pt}
{\fontsize{14}{18}\selectfont A canonical Prikry layer at the ultraexacting boundary\par}
\vspace{14pt}
{\color{Muted}Prepared for Vladimir Reshetnikov\hfill 18 September 2026\par}
\vspace{3pt}
{\color{Sage}\hrule height .6pt}
\vspace{12pt}

\textbf{Abstract.}
Let $\lambda$ be ultraexacting. We construct a class $q$ of cofinal
$\omega$-subsets of $\lambda$ modulo finite difference with the following
property. For every parameter $p\in V_\lambda$, the inner model
\[
 M_{q,p}=\{x:\tc(\{x\})\subseteq\OD(q,p)\}
\]
contains a normal measure on $\lambda$, defined externally by eventual
membership along $q$. Every representative of $q$ is Prikry generic over
$M_{q,p}$ for this measure. The same $q$ works for all such $p$.
The key new step is a regressive-function argument at the \emph{first}
critical point, combined with the archive's minimal-parameter-rank method.

\begin{assessment}{The resulting consistency statements}
\textbf{A definability obstruction.} Ultraexactingness contradicts the
principle that every cofinal tail class admits a splitting subset of
$\lambda$ ordinal definable from that class and a parameter in $V_\lambda$.

\textbf{An enhanced equiconsistency package.} The theory asserting an
ultraexacting $\lambda$ with $V_\lambda\subseteq\HOD$, together with the
canonical normal measure and Prikry layer constructed here, is
equiconsistent with $\ZFC+\Izero$. The underlying
ultraexacting--$\Izero$ comparison and the coding step are imported;
the additional canonical structure is proved here.

\textbf{A stationary-set obstruction.} A strongly compact cardinal below
$\lambda$ prevents the canonical Prikry layer from being the entire
universe. Any further outer extension satisfying that hypothesis must
collapse its successor of $\lambda$ or destroy a specific stationary set.
The stationary mechanism is the archive's earlier one, now applied to the
newly constructed layer.
\end{assessment}

\textbf{Scope and status.}
This is a conventional, unrefereed research argument, not a claim of an
unconditional inconsistency of ultraexacting cardinals or $\Izero$.
All new implications are proved in English. No new Lean verification is
claimed. The measure is complete \emph{inside the indicated inner model},
not in the ambient universe, where $\cf(\lambda)=\omega$.
The contribution is a continuation of the supplied archive; worldwide
priority is not asserted.

\newpage
\tableofcontents

\newpage
\section{What is added to the archive}\label{sec:delta}

The supplied \nolinkurl{Cardinals3.zip} contains the second-edition
\emph{Large Cardinals Synthesis}, the \emph{Unified Report}, and Lean
sources. The synthesis already proves the completeness barriers,
sharp-width theorems, and rigidity of certain classes in
$D_\lambda/{=^*}$; it also contains an $\Izero$-calibrated coding package
and a successor-stationarity obstruction \cite{synthesis,plan}.
Those results are not being counted again as new results here.

The decisive change of question is from
\emph{``Can a tail class define a small family of representatives?''}
to
\emph{``Which subsets of $\lambda$ contain almost all of a critical
sequence, and do they form a normal measure in an inner model?''}
Rigidity alone gives no answer to the normality question.
Indeed, Section~\ref{sec:sharp} gives an explicit recoding which retains
the same relative-definability model but makes the corresponding tail
ultrafilter nonnormal.

\begin{center}
\small
\renewcommand{\arraystretch}{1.19}
\begin{tabularx}{\textwidth}{@{}P{.27\textwidth}YY@{}}
\toprule
\textbf{Issue}&\textbf{Archive baseline}&\textbf{Extension here}\\\midrule
Critical-sequence class&Fixed modulo finite difference; relative rigidity&No relative-definable splitter; every relative-definable regressive function is eventually constant\\
Inner model at $\lambda$&Regularity, and inaccessibility in the coded boundary case&An actual \emph{internal normal measure} in each fixed-parameter relative HOD\\
Prikry phenomenon&An iteration-based consistency construction&An intrinsic Prikry layer inside \emph{every} ultraexacting universe, with no iterability assumption added\\
Consistency strength&Ultraexactingness calibrated by $\Izero$&The canonical measure/layer and full low-rank agreement can be added without increasing that calibration\\
Successor obstruction&Collapse or destruction for an arbitrary regularizing inner model&A concrete canonical intermediate model and a named stationary witness\\
\bottomrule
\end{tabularx}
\end{center}

The phenomenon that critical sequences can be Prikry generic over an
iteration model is classical. A related iteration-based use occurs in
\cite[Section~5.2]{abgl}. Our assertion concerns instead a
\emph{specific hereditary ordinal-definability class formed in the
ambient universe from the tail class itself}. It includes uniformity over
all additional parameters of rank below $\lambda$. Neither a well-founded
iteration of the given witness nor an identification of the universe as
a forcing extension is assumed.

\subsection{Dependency boundary}

We use the local Kunen inconsistency theorem: in $\ZFC$ there is no
nontrivial elementary self-embedding of $V_{\mu+2}$
\cite{kunen}. For the final equiconsistency assertion we additionally use
two results of Aguilera--Bagaria--Goldberg--L\"ucke: the equiconsistency of
ultraexactingness with $\Izero$, and forcing which preserves an
ultraexacting $\lambda$ while making $V_\lambda\subseteq\HOD$
\cite[Theorem~A and Proposition~3.9]{abgl}.
The normal-measure construction itself does not use the consistency
comparison or the coding theorem.

The genericity criterion is Mathias's criterion \cite{mathias};
Section~\ref{sec:prikryproof} supplies a proof, including the required
Prikry lemmas. The strong-compactness application also includes its
reflection argument. Thus those applications do not hide a new
unproved forcing or reflection lemma.

\section{Parameters, tail classes, and inner models}\label{sec:notation}

Work in an ambient universe $V\models\ZFC$. All unadorned ranks,
cardinals, and definability classes are computed in $V$.
For a set $x$, $\tc(\{x\})$ is the transitive closure of the singleton
$\{x\}$, so in particular it includes $x$.

\begin{definition}[Ultraexactingness]\label{def:uex}
An uncountable cardinal $\lambda$ is \emph{ultraexacting} if, for every
sufficiently large cardinal $\zeta>\lambda$, there are
\[
 X\prec V_\zeta,\qquad
 V_\lambda\cup\{\lambda\}\subseteq X,\qquad
 j:X\longrightarrow V_\zeta
\]
such that $j$ is elementary, $j(\lambda)=\lambda$,
$j\restr\lambda\ne\id$, and $j\restr V_\lambda\in X$.
We use the equivalent formulation with arbitrarily large heights; the
usual definition quantifies over every height above $\lambda$
\cite[Definition~3.3 and Lemma~3.2]{abl}.
\end{definition}

The domain $X$ is not assumed transitive. In particular,
$x\in X$ does not imply $x\subseteq X$. Only the explicit inclusion
$V_\lambda\subseteq X$ and elementary definable closure will be used.

\begin{definition}[Tail quotient]\label{def:tail}
For $\cf(\lambda)=\omega$, set
\[
 D_\lambda=\{a\subseteq\lambda:\ot(a)=\omega,\ \sup a=\lambda\},
 \qquad Q_\lambda=D_\lambda/{=^*},
\]
where $a=^*b$ means that $a\triangle b$ is finite. Write
$[a]=\{b\in D_\lambda:b=^*a\}$. For $q\in Q_\lambda$ and $A\subseteq\lambda$,
put
\[
 \Tail(q,A)\quad\Longleftrightarrow\quad
 (\exists a\in q)\ |a\setminus A|<\omega.
\]
The existential quantifier may equivalently be universal. Say that
$A$ \emph{splits} $q$ if both $a\cap A$ and $a\setminus A$ are infinite
for one, equivalently every, $a\in q$.
\end{definition}

Every representative has rank $\lambda$, and the full class $q$ is an
element of $V_{\lambda+2}$. Passing from a representative to its class
forgets finitely many entries; it does not supply a preferred
representative.

\begin{definition}[The parameter convention]\label{def:hod}
$\OD(q,p)$ consists of the sets definable in $V$ from ordinals and the
\emph{two individual set parameters} $q,p$. Equivalently one can use a
unique definition in a sufficiently large rank segment with these
parameters. Set
\[
 \mathcal O_q=\bigcup_{p\in V_\lambda}\OD(q,p),
 \qquad
 M_{q,p}=\{x:\tc(\{x\})\subseteq\OD(q,p)\},
 \qquad M_q=M_{q,\varnothing}.
\]
A finite tuple of members of $V_\lambda$ can be coded by one member of
$V_\lambda$, since $\lambda$ is a limit cardinal.
\end{definition}

\begin{assessment}{A parameter is not a family of permitted parameters}
In $\OD(q,p)$ the class $q$ is \emph{one set parameter}.
Its members are not individually available as parameters. Also,
$M_{q,p}$ is not $L(q,p)$, and it need not contain $q$ or $p$.
If $q$ were an element of this transitive inner model, all its
representatives would be elements too. The construction below proves
precisely that this does not happen.
\end{assessment}

\begin{lemma}[Fixed-parameter hereditary OD]\label{lem:hod}
For arbitrary sets $q,p$, $M_{q,p}$ is a definable transitive inner model
of $\ZFC$. The class $\OD(q,p)$ has a $(q,p)$-definable set-like
well-order, and every nonempty definable subclass of it has a least
member in that order.
\end{lemma}

\begin{proof}
Code a definition by a rank height, a formula number, and a finite tuple
of ordinal parameters. The rank height is required to contain $q,p$
and all the other parameters. A code is valid when its formula uniquely
defines an object in that rank segment. Order codes first by a bounding
ordinal for all their entries, and then by a fixed well-order of the
set of codes with that bound. Every initial segment is a set, and every
nonempty class of valid codes has a least code. Assign each definable
object its least code. This gives the required definable set-like
well-order. The satisfaction relation used here is the definable
satisfaction relation of a \emph{set} structure, not a truth predicate
for $V$.

Write $M=M_{q,p}$. Its definition immediately gives transitivity and
membership of every ordinal. Whenever an object is definable from
ordinals, $q,p$, and finitely many members of $\OD(q,p)$, replace the
latter parameters by their definitions. The object is again in
$\OD(q,p)$. We use this substitution observation repeatedly.

For $x,y\in M$, pairing and union are ordinal definable from $q,p$ and
have only hereditarily ordinal-definable members, so their usual values
belong to $M$. The same argument gives Infinity. For $x\in M$, the set
\[
 \{y\subseteq x:y\in M\}
\]
exists in $V$, is definable from $x,q,p$, and has all its members in $M$.
It therefore belongs to $M$ and is its internal power set of $x$.
For any fixed formula, relativizing its quantifiers to the definable
class $M$ is an ordinary formula of $V$. Separation in $V$ and the
substitution observation give Separation in $M$. If such a formula
defines a function on a set of $M$ with values in $M$, Replacement in
$V$ produces the range; its definition and its hereditary constituents
put the range in $M$. Extensionality and Foundation follow by
transitivity.

Finally, restrict the canonical $\OD(q,p)$ well-order to a set $x\in M$.
This restriction is definable from $q,p$ and $x$, hence from $q,p$ and
ordinals. Its transitive closure consists of the relation itself,
finite set constructions from elements of $x$, and their transitive
closures. All these are in $\OD(q,p)$. The restriction thus belongs to
$M$ and well-orders $x$ there. This proves Choice.
\end{proof}

We will not infer Choice for a larger hereditary class obtained by
allowing \emph{all} of $V_\lambda$ as separate parameters at once.
The fixed-parameter models in Lemma~\ref{lem:hod} suffice, and avoid
that extra issue entirely.

\section{One witness, uniformly correct for all parameters}\label{sec:witness}

\subsection{The critical sequence}

\begin{lemma}[Critical-sequence normalization]\label{lem:critical}
Let $j:X\to V_\zeta$ witness ultraexactingness at a height
$\zeta>\lambda+\omega$, and put
\[
 e=j\restr V_\lambda,\quad \kappa=\crit(e),\quad
 \kappa_n=e^n(\kappa),\quad c=\{\kappa_n:n<\omega\}.
\]
Then $e:V_\lambda\to V_\lambda$ is elementary,
\[
 \sup_{n<\omega}\kappa_n=\lambda,\qquad
 e(\xi)>\xi\quad(\kappa\le\xi<\lambda),\qquad
 e\restr V_\kappa=\id.
\]
Moreover $c,q=[c]\in X$, and
\[
 j(c)=c\setminus\{\kappa\},\qquad j(q)=q.
\]
\end{lemma}

\begin{proof}
The object $V_\lambda$ is definable from $\lambda$ at the chosen height,
so belongs to $X$ and is fixed by $j$. Relativizing any formula to
$V_\lambda$, and using $V_\lambda\subseteq X$, shows that $e$ is
an elementary self-embedding. This reasoning does not require
transitivity of $X$.

The critical point is below $\lambda$ because $j\restr\lambda$ is
nonidentity. Elementary embeddings are strictly increasing on
ordinals and never move an ordinal downward. It follows that
$\kappa<\kappa_1<\kappa_2<\cdots<\lambda$.
Suppose $\mu=\sup_n\kappa_n<\lambda$. The sequence then belongs to
$V_\lambda$, and elementarity shows $e(\mu)=\mu$. The restriction to
$V_{\mu+2}$ is a nonidentity elementary self-embedding, contrary to
local Kunen. Here $\mu+2<\lambda$: $\lambda$ is an infinite cardinal,
whereas $\mu<\lambda$. Thus $\sup_n\kappa_n=\lambda$, and in particular
$\lambda$ is a limit cardinal of cofinality $\omega$.
For $\kappa_n\le\xi<\kappa_{n+1}$ we have
$e(\xi)\ge\kappa_{n+1}>\xi$.

An induction on rank proves $e\restr V_\kappa=\id$. At a rank
$\alpha<\kappa$, the ordinal $\alpha$ and $V_\alpha$ are fixed. If all
members of $V_\alpha$ are already fixed, membership in any
$x\subseteq V_\alpha$ is preserved pointwise, and $e(x)$ is still a
subset of $V_\alpha$. Hence $e(x)=x$.

Since the \emph{set} $e$ belongs to $X$, its first moved ordinal and its
$\omega$-iteration on that ordinal are definable in $V_\zeta$ from $e$.
They belong to $X$ by elementary definable closure. In particular
$c\in X$. The increasing enumeration of $c$ is in $X$ as well.
$j$ fixes every natural number, so applying $j$ to that enumeration
shows $j(c)=j``c=\{\kappa_{n+1}:n<\omega\}$.

Finally $q$ is definable from $c,\lambda$. All subsets of $\lambda$
and their finite-modification classes occur below the chosen height,
so this definition computes the actual class, not a truncated one.
Elementarity gives $j(q)=[j(c)]=[c]=q$.
\end{proof}

\subsection{Uniform reflection is chosen before the witness}

The argument must handle \emph{all} ordinal definitions, including
ones using ordinals above the critical point. It would be incorrect
to choose an arbitrary definition of a counterexample and simply
assert that its ordinal parameters are fixed.

For each of the finitely many counterexample properties used below,
fix ordinary first-order formulas expressing membership in
$\OD(q,p)$, existence of a counterexample, its least permitted
parameter rank, and selection of the canonical least counterexample.
These are fixed formulas, even though they quantify over codes for
arbitrarily many definitions.

By the reflection theorem choose, \emph{before choosing $j$}, a cardinal
$\zeta>\lambda+\omega$ such that $V_\zeta$ is correct for these formulas
and their required subformulas, for \emph{every} tuple of parameters
in $V_\zeta$. One may choose a sufficiently large finite $n$ and take
$\zeta$ in $C^{(n)}$ and in the club of cardinal closure points.
This is finite-level correctness, not the assertion
$V_\zeta\prec V$.

In this report the two counterexample predicates concern subsets of
$\lambda$ and functions $\lambda\to\lambda$. Such objects are all in
$V_{\lambda+\omega}$; therefore the selected counterexamples lie
below $\zeta$. Choose the ultraexacting witness at this $\zeta$, and
form $c,q$ as in Lemma~\ref{lem:critical}.

\begin{lemma}[Fixed-counterexample principle]\label{lem:minrank}
Fix one of the preceding counterexample formulas $B(x,q,\lambda)$.
If
\[
 (\exists p\in V_\lambda)(\exists x\in\OD(q,p))\ B(x,q,\lambda),
\]
then there is such a counterexample $x\in X$ with $j(x)=x$.
\end{lemma}

\begin{proof}
Let $\rho<\lambda$ be least such that some $p\in V_\rho$ permits a
counterexample in $\OD(q,p)$. The minimum exists by the displayed
hypothesis and the fact that $\lambda$ is a limit ordinal.
Its defining formula uses only $q,\lambda$. Uniform correctness
ensures that $V$ and $V_\zeta$ compute the same $\rho$. Since
$q,\lambda\in X$ are fixed, elementary definable closure puts
$\rho$ in $X$, and elementarity gives $j(\rho)=\rho$.
By Lemma~\ref{lem:critical}, $\rho<\kappa$.

Choose any witnessing $p\in V_\rho$. This $p$ belongs to
$V_\kappa\subseteq X$ and is fixed by $j$. Among the counterexamples
in $\OD(q,p)$ choose the least under the canonical well-order of
Lemma~\ref{lem:hod}. Its selection is a unique definition from
$q,p,\lambda$. The height was chosen correct for that selection
formula, so the selected $x$ is the same in $V$ and $V_\zeta$.
Elementarity puts it in $X$ and fixes it.

Notice that no ordinal parameter in the \emph{original} definition of
$x$ was assumed fixed. The minimization and canonical selection
replace that invalid assumption by a unique fixed definition.
\end{proof}

The minimal-rank mechanism is inherited from the synthesis's
cofinal-orbit-rigidity argument \cite{synthesis}. Its use for the two
new counterexample properties is what permits a single tail class to
work uniformly for all low-rank parameters.

\section{The two tail theorems}\label{sec:tailtheorems}

Fix from now on the uniformly correct witness of
Section~\ref{sec:witness}, its critical sequence $c$, and $q=[c]$.

\begin{theorem}[No relative-definable splitter]\label{thm:nosplit}
For every $A\subseteq\lambda$ with $A\in\mathcal O_q$, exactly one of
\[
 \Tail(q,A),\qquad \Tail(q,\lambda\setminus A)
\]
holds. Equivalently, every such $A$ either contains almost all of $c$
or meets $c$ in only finitely many points.
\end{theorem}

\begin{proof}
If neither alternative holds, $A$ splits $q$. Splitting is expressible
using $q,\lambda,A$, without selecting a representative. Apply
Lemma~\ref{lem:minrank} to the property of being a splitting subset of
$\lambda$. It gives a splitting $A\in X$ with $j(A)=A$.
For every $n<\omega$, elementarity of $j$, and of the inclusion
$X\prec V_\zeta$, gives
\[
 \kappa_n\in A
 \quad\Longleftrightarrow\quad
 j(\kappa_n)\in j(A)
 \quad\Longleftrightarrow\quad
 \kappa_{n+1}\in A.
\]
Thus the membership pattern on $c$ is constant, contradicting that
$A$ splits $q$. Both alternatives cannot hold because $c$ is infinite.
\end{proof}

\begin{definition}[Eventual regressivity]
A function $f:\lambda\to\lambda$ is \emph{eventually regressive on
$q$} if $f(\alpha)<\alpha$ for all but finitely many $\alpha$ in a
representative of $q$. It is \emph{eventually constant on $q$} if some
$\xi<\lambda$ satisfies $\Tail(q,\{\alpha:f(\alpha)=\xi\})$.
Both definitions are independent of the representative.
\end{definition}

\begin{theorem}[Tail normality]\label{thm:tailnormal}
Every $f:\lambda\to\lambda$ in $\mathcal O_q$ that is eventually
regressive on $q$ is eventually constant on $q$.
\end{theorem}

\begin{proof}
Suppose the theorem fails. Its failure is a first-order property
$B(f,q,\lambda)$: $f$ is a total function of the indicated type, is
eventually regressive, and is not eventually constant.
Lemma~\ref{lem:minrank} supplies a counterexample $f\in X$ fixed by $j$.

For every $n$, $f(\kappa_n)$ belongs to $X$, because it is uniquely
definable there from $f,\kappa_n$. It is below $\lambda$, so $j$ agrees
with $e$ at that value. Consequently
\begin{equation}\label{eq:recurrence}
 f(\kappa_{n+1})=e(f(\kappa_n))\qquad(n<\omega).
\end{equation}
Moreover, elementarity gives the two-way equivalence
\[
 f(\kappa_n)<\kappa_n
 \quad\Longleftrightarrow\quad
 f(\kappa_{n+1})<\kappa_{n+1}.
\]
Eventual regressivity therefore implies regressivity at \emph{every}
entry, in particular
\[
 \xi:=f(\kappa_0)<\kappa_0=\kappa.
\]
This value is fixed by $e$. Equation~\eqref{eq:recurrence}, by ordinary
induction on $n$, now yields $f(\kappa_n)=\xi$ for every $n$.
This contradicts failure of eventual constancy.
\end{proof}

\begin{assessment}{Where normality enters}
For subsets, shift invariance only says that a binary membership
pattern is constant. For regressive functions, the first value lies
\emph{below the critical point}, hence is fixed. That extra fact forces
all subsequent values to be the same. Starting with an arbitrary
orbit above the critical point does not supply this fact.
\end{assessment}

\begin{corollary}[No definable representative]\label{cor:norep}
No member of $q$ belongs to $\mathcal O_q$. In particular,
$q\notin M_{q,p}$ and $q\cap M_{q,p}=\varnothing$ for all $p\in V_\lambda$.
\end{corollary}

\begin{proof}
If $a\in q\cap\OD(q,p)$, let $A$ consist of the even-indexed entries of
its increasing enumeration. Then $A\in\OD(q,p)$ and $A$ splits $q$,
contradicting Theorem~\ref{thm:nosplit}. If $q\in M_{q,p}$, transitivity
would place every member of $q$ in $M_{q,p}\subseteq\OD(q,p)$.
\end{proof}

\section{The internal normal-measure theorem}\label{sec:measure}

\begin{theorem}[Canonical internal normal measures]\label{thm:measure}
Let $\lambda$ be ultraexacting. There is $q\in Q_\lambda$ such that, for
\emph{every} $p\in V_\lambda$, the set
\begin{equation}\label{eq:U}
 U_{q,p}=\{A\in\Pow(\lambda)\cap M_{q,p}:\Tail(q,A)\}
\end{equation}
belongs to $M_{q,p}$ and is a nonprincipal normal
$\lambda$-complete ultrafilter on $\lambda$ \emph{in that model}.
In particular
\[
 M_{q,p}\models\text{``$\lambda$ is measurable.''}
\]
The same $q$ can be used simultaneously for all $p\in V_\lambda$.
\end{theorem}

\begin{proof}
Take the $q$ from Section~\ref{sec:tailtheorems}. Fix $p\in V_\lambda$,
and abbreviate $M=M_{q,p}$ and $U=U_{q,p}$.

\textbf{The measure is an element of the model.}
The right side of \eqref{eq:U} defines a set in $V$ from $q,p$ and the
ordinal $\lambda$, using the definable predicate for $M$.
Hence $U\in\OD(q,p)$. Every member $A$ of $U$ is a subset of an ordinal
and belongs to $M$. Its transitive closure is contained in
$\OD(q,p)$. Together with $U\in\OD(q,p)$, this proves
$\tc(\{U\})\subseteq\OD(q,p)$, or $U\in M$.
This is stronger than merely having an external ultrafilter on the
Boolean algebra $\Pow(\lambda)\cap M$.

\textbf{Ultrafilter and nonprincipality.}
The set $\lambda$ belongs to $U$ and the empty set does not. If two sets
contain almost all of a representative of $q$, their intersection does
too. Upward closure is immediate. For every $A\in\Pow(\lambda)\cap M$,
Theorem~\ref{thm:nosplit} selects exactly one of $A$ and
$\lambda\setminus A$. These properties are the ultrafilter axioms for
$\Pow^M(\lambda)$. Every final segment of $\lambda$ is in $U$, since
representatives are increasing cofinal $\omega$-sets. No bounded set,
and in particular no singleton, is in $U$.

\textbf{Normality.}
Let $S\in U$ and let $f\in M$ be a regressive function on $S\setminus\{0\}$
with values below $\lambda$. Extend it to a total function
$\bar f:\lambda\to\lambda$ by setting it equal to zero outside
$S\setminus\{0\}$. This total function is in $M$ and hence in
$\OD(q,p)$. Since a representative of $q$ is eventually in
$S\setminus\{0\}$, $\bar f$ is eventually regressive on $q$.
Theorem~\ref{thm:tailnormal} supplies $\xi<\lambda$ with
$\Tail(q,\{\alpha:\bar f(\alpha)=\xi\})$.
Intersecting with $S\setminus\{0\}$ gives
\[
 \{\alpha\in S\setminus\{0\}:f(\alpha)=\xi\}\in U.
\]
The fiber is in $M$ by Separation. This proves the internal normality
statement, with all its quantifiers interpreted in $M$.

\textbf{$\lambda$-completeness.}
Let $\tau<\lambda$ and let
$\langle A_i:i<\tau\rangle\in M$ be a sequence of members of $U$.
For $\tau=0$ the assertion is immediate. Otherwise suppose
$A=\bigcap_{i<\tau}A_i\notin U$. Then
$S=\lambda\setminus A\in U$. In $M$ define
\[
 f(\alpha)=\min\{i<\tau:\alpha\notin A_i\}
 \qquad(\alpha\in S).
\]
On $S\setminus(\tau+1)\in U$ this is regressive. Normality yields an
$i<\tau$ such that
\[
 T=\{\alpha\in S\setminus(\tau+1):f(\alpha)=i\}\in U.
\]
But $T\cap A_i=\varnothing$, contradicting the filter property.
Thus the intersection belongs to $U$.

Finally $\lambda$ remains an uncountable cardinal in $M$, since an
inner-model bijection contradicting that assertion would be a
bijection in $V$ as well. The preceding properties therefore witness
its measurability in $M$. Uniformity in $p$ follows because the two tail
theorems were proved for all of $\mathcal O_q$ using the same $q$.
\end{proof}

\begin{corollary}[Regularity and strong limit, internally]\label{cor:inacc}
Each $M_{q,p}$ regards $\lambda$ as strongly inaccessible.
\end{corollary}

\begin{proof}
We give the standard deductions to make the internal quantifiers
explicit. Work entirely in $M=M_{q,p}$. A cofinal sequence
$\langle\alpha_i:i<\tau\rangle$ with $\tau<\lambda$ would give the
$\tau$ many members $\lambda\setminus(\alpha_i+1)$ of $U$ with empty
intersection, contradicting completeness. Hence $\lambda$ is regular.

Suppose $\eta<\lambda$ and $2^\eta\ge\lambda$. By Choice in $M$ there
is an injection $h:\lambda\to\Pow(\eta)$. For each $\xi<\eta$, let
$B_\xi=\{\alpha<\lambda:\xi\in h(\alpha)\}$. Select the one of
$B_\xi$ and $\lambda\setminus B_\xi$ belonging to $U$.
Their intersection is in $U$ by completeness, and all its members
have exactly the same $h$-value. Injectivity would make that
intersection have at most one member, contradicting nonprincipality.
Thus $2^\eta<\lambda$ for all $\eta<\lambda$.
\end{proof}

\begin{proposition}[Uniqueness and coherence]\label{prop:coherence}
For the $q$ of Theorem~\ref{thm:measure}:
\begin{enumerate}
\item $U_{q,p}$ is the unique ultrafilter on $\Pow^{M_{q,p}}(\lambda)$
whose every member contains almost all of a representative of $q$.
\item If $\OD(q,p)\subseteq\OD(q,r)$, then $M_{q,p}\subseteq M_{q,r}$ and
\[
 U_{q,p}=U_{q,r}\cap\Pow^{M_{q,p}}(\lambda).
\]
\item Given $p_0,p_1\in V_\lambda$, coding the pair by
$r\in V_\lambda$ gives a common larger model $M_{q,r}$ and coherent
measures as in part~(2).
\end{enumerate}
\end{proposition}

\begin{proof}
For part (1), the assumed tail property makes every member of the
other ultrafilter a member of $U_{q,p}$. Two ultrafilters on the same
Boolean algebra cannot be properly included one in the other: the
complement of any proposed missing member gives a contradiction.
For part (2), the hereditary definitions give the model inclusion,
and the same formula $\Tail(q,A)$ decides both measures on their
common subsets. Coding gives part (3).
\end{proof}

\subsection{Why ambient countable incompleteness is harmless}

Let $c=\{\kappa_n:n<\omega\}$ be the chosen critical sequence.
Every set $T_n=\lambda\setminus(\kappa_n+1)$ is in $U_{q,p}$,
yet in the ambient universe
\begin{equation}\label{eq:externalincomplete}
 \bigcap_{n<\omega}T_n=\varnothing.
\end{equation}
This explicitly shows that $U_{q,p}$ is not externally countably
complete. The sequence $\langle T_n:n<\omega\rangle$ is not an element
of $M_{q,p}$: otherwise its internal completeness would contradict
\eqref{eq:externalincomplete}. Each individual $T_n$ is in the model,
but the external enumeration of them is not.

Nor does \eqref{eq:externalincomplete}, by itself, prove that the
ultrapower of $M_{q,p}$ using only its \emph{internal} functions is
externally ill-founded. An external sequence of sets and an internal
sequence of functions are different data. We make no iterability or
external ultrapower well-foundedness assertion.

\section{Canonical Prikry layers}\label{sec:layers}

\begin{definition}[Prikry forcing and the filter from a sequence]
Let $M$ be a transitive inner model of $\ZFC$ containing a normal measure
$U$ on its measurable cardinal $\lambda$. A condition in $\PP_U^M$ is
$(s,A)$, where $s$ is a finite increasing sequence of ordinals below
$\lambda$, $A\in U$, and every member of $A$ is above every entry of $s$.
The stronger-condition convention is used:
$(t,B)\le(s,A)$ if $t$ end-extends $s$, its new entries belong to $A$,
and $B\subseteq A$. A direct extension, written $\le^*$, keeps the stem.

For $a\in D_\lambda$, identify $a$ with its increasing enumeration and
put
\[
 G_a=\{(s,A)\in\PP_U^M:s\text{ is an initial segment of }a,
              \ a\setminus\ran(s)\subseteq A\}.
\]
\end{definition}

\begin{lemma}[Mathias's criterion]\label{lem:mathias}
If $a\in D_\lambda$ satisfies $a\subseteq^* A$ for every $A\in U$, then
$G_a$ is $\PP_U^M$-generic over $M$.
\end{lemma}

A proof of this classical criterion, entirely with the internal
sequences and dense sets of $M$, appears in
Section~\ref{sec:prikryproof}. This is the form needed here; no
countability assumption on $M$ is made.

\begin{theorem}[The canonical layer]\label{thm:layer}
Let $q$ be as in Theorem~\ref{thm:measure}. For every $p\in V_\lambda$ and
every $a\in q$, $a$ is Prikry generic over $M_{q,p}$ for $U_{q,p}$.
Define
\[
 K_{q,p}=M_{q,p}[G_a]=M_{q,p}[a].
\]
Then $K_{q,p}$ is independent of the choice of $a\in q$, and
\[
 M_{q,p}\subsetneq K_{q,p}\subseteq V,\qquad q\in K_{q,p}.
\]
The two inner models have the same cardinals and
\begin{equation}\label{eq:rankagreement}
 V_\lambda^{M_{q,p}}=V_\lambda^{K_{q,p}},
 \qquad
 \cf^{K_{q,p}}(\lambda)=\omega,
 \qquad
 (\lambda^+)^{M_{q,p}}=(\lambda^+)^{K_{q,p}}.
\end{equation}
\end{theorem}

\begin{proof}
By the definition of $U_{q,p}$, every $A\in U_{q,p}$ contains almost
all of \emph{every} representative $a\in q$.
Lemma~\ref{lem:mathias} therefore applies. The resulting generic filter
is a set of $V$, and evaluating names from $M_{q,p}$ produces a
transitive inner model of $V$. It is a genuine set-forcing extension
of $M_{q,p}$, and hence satisfies $\ZFC$.
The generic filter is definable from $a$ and the forcing parameters,
while $a$ is recovered as the union of its stems. This explains
$M_{q,p}[G_a]=M_{q,p}[a]$ without treating an arbitrary notation
$M[a]$ as an independently given forcing extension.

If $a,b\in q$, the finite set $a\triangle b$ belongs to $M_{q,p}$,
because every finite set of ordinals does. Thus $a$ and $b$ are mutually
recoverable over $M_{q,p}$ and give the same extension. This proves
that $K_{q,p}$ is well-defined. The full equivalence class $q$ is in
$K_{q,p}$: it consists precisely of the finite modifications of $a$,
and all finite subsets of $\lambda$ are already present. This
calculation gives the actual ambient class, not merely a smaller
class of representatives computed internally. Properness follows
from Corollary~\ref{cor:norep}.

The standard preservation facts for normal Prikry forcing are proved
in Section~\ref{sec:prikryproof}: there are no new bounded subsets of
$\lambda$, the forcing is $\lambda^+$-cc internally, and all cardinals
are preserved. The sequence $a$ changes the cofinality of $\lambda$
to $\omega$. These give the successor assertion.

For rank agreement, induct on $\alpha<\lambda$. In $M_{q,p}$,
$\lambda$ is strongly inaccessible by Corollary~\ref{cor:inacc}, so
$|V_\alpha|<\lambda$. If the rank segments up to $\alpha$ agree, a new
subset of $V_\alpha$ could be coded, using a bijection in $M_{q,p}$,
by a new bounded subset of $\lambda$. There are no such subsets.
Successor stages therefore agree, and limit stages are unions.
This proves \eqref{eq:rankagreement}.
\end{proof}

\begin{remark}[A layer is not an ambient ground]
The theorem proves that $M_{q,p}$ is a ground of $K_{q,p}$. It does not
prove that either model is a ground of $V$, or that $K_{q,p}=V$.
No quotient forcing from $K_{q,p}$ to $V$ is supplied. This distinction
is essential in the stationary-set application below.
\end{remark}

\begin{corollary}[The coded boundary]\label{cor:boundary}
Assume additionally $V_\lambda\subseteq\HOD^V$. Then
$M_{q,p}=M_q$ and $U_{q,p}=U_q$ for every $p\in V_\lambda$.
The common Prikry layer $K_q$ satisfies
\[
 V_\lambda^{M_q}=V_\lambda^{K_q}=V_\lambda^V,
 \qquad
 M_q\models\text{``$\lambda$ is measurable.''}
\]
Nevertheless no representative of $q$ belongs to $M_q$.
\end{corollary}

\begin{proof}
Each $p\in V_\lambda$ is ordinal definable in $V$, so it can be
eliminated from any $\OD(q,p)$ definition by substitution. Hence
$\OD(q,p)=\OD(q)$ and the hereditary models are equal. Every element
of $V_\lambda$ and its transitive closure are in $\HOD\subseteq M_q$,
so $V_\lambda\subseteq M_q$. Transitivity and inclusion in $V$ give the
rank equalities. The remaining statements were proved above.
\end{proof}

The measurable cardinal here is $\lambda$ in $M_q$, not in ordinary
$\HOD$ and not in $V$. There is no argument eliminating the high-rank
parameter $q$ from the definition of $U_q$.

\section{Sharpness and diagnostic counterexamples}\label{sec:sharp}

\subsection{The low-rank parameter bound cannot be raised}

Allow the additional parameter $p=c$, which belongs to
$V_{\lambda+1}$ but not $V_\lambda$. The subset
\[
 A=\{\kappa_{2n}:n<\omega\}
\]
is definable from $c$ and splits $q$. Also the function which assigns
to $\alpha$ the largest member of $c\cap\alpha$ when that set is
nonempty, and assigns zero otherwise, is definable from $c$.
At $\kappa_n$ for $n>0$ it takes the value $\kappa_{n-1}$.
It is eventually regressive and not eventually constant.
Thus both tail theorems genuinely fail when parameters from
$V_{\lambda+1}$ are allowed indiscriminately.

\subsection{A fixed recoding with a nonnormal tail ultrafilter}

Let
\[
 c^+=\{\kappa_n+1:n<\omega\},\qquad q^+=[c^+].
\]
The ordinal map $h(\alpha)=\alpha+1$ sends $\lambda$ into $\lambda$.
The class $q^+$ is definable from $q$, without selecting a
representative: apply $h$ to any representative and take its
finite-difference class. Conversely, apply predecessor to the
successor ordinals in a representative of $q^+$ and ignore the
finitely many nonsuccessors. The resulting class is $q$.
Both operations are well-defined modulo finite difference.

Consequently
\begin{equation}\label{eq:recodehod}
 \OD(q,p)=\OD(q^+,p),\qquad M_{q,p}=M_{q^+,p}.
\end{equation}
Also $q^+\in X$ and $j(q^+)=q^+$, since the recoding is definable from
fixed parameters. The no-splitting property transfers: if $A$ split
$q^+$, then $h^{-1}[A]$ would split $q$ and would be definable from the
same parameters. Thus the tail rule on $q^+$ is an ultrafilter on the
same internal Boolean algebra.

\begin{proposition}[Normality distinguishes the critical class]\label{prop:shift}
In the common model $M=M_{q,p}=M_{q^+,p}$, the tail ultrafilter
\[
 U^+=\{A\in\Pow^M(\lambda):\Tail(q^+,A)\}
\]
is $\lambda$-complete and nonprincipal but is not normal.
\end{proposition}

\begin{proof}
The recoding gives
\[
 A\in U^+\quad\Longleftrightarrow\quad h^{-1}[A]\in U_{q,p}.
\]
Preimages commute with intersections, so internal
$\lambda$-completeness follows. The formula also shows that $U^+$ is
an element of $M$; alternatively use its external definition and the
same hereditary argument as for $U_{q,p}$.
The set $S$ of successor ordinals below $\lambda$ belongs to $U^+$.
On $S$ the function $f(\alpha+1)=\alpha$ is an ordinal-definable
regressive function. Every fiber is a singleton, hence is outside
$U^+$. This contradicts the normality condition.
\end{proof}

This example pinpoints a possible false strengthening: neither
$j(q)=q$ nor equality of the relative HOD models forces normality of
\emph{every} tail presentation. In the proof of
Theorem~\ref{thm:tailnormal}, the inequality needed is
$f(\kappa)<\kappa$, not merely $f(\beta)<\beta$ at an arbitrary orbit
root $\beta>\kappa$.

\subsection{What no-splitting does not assert}

Choice still supplies splitting subsets for all $q\in Q_\lambda$.
Our obstruction concerns their restricted definability, not their
existence as sets. The tail measure need not decide arbitrary ambient
subsets of $\lambda$: the even and odd pieces of $c$ are already
counterexamples. Also the normal measures belonging to different
fixed-parameter models need not form a single internally indexed
sequence in any one of them. The coherence of
Proposition~\ref{prop:coherence} is not a claim of such internal
uniform enumeration.

\section{Conditional inconsistency and equiconsistency}\label{sec:consistency}

\subsection{A splitting principle incompatible with ultraexactingness}

\begin{definition}[Relative-definable splitting]
For an uncountable cardinal $\lambda$ of cofinality $\omega$, let
$\NS(\lambda)$ be the assertion
\[
 \forall q\in Q_\lambda\ \exists p\in V_\lambda\ \exists A\subseteq\lambda
 \ \bigl(A\in\OD(q,p)\ \wedge\ A\text{ splits }q\bigr).
\]
This is a pointwise definability principle; it does not demand a
uniform choice of the definitions or their parameters.
\end{definition}

\begin{theorem}[Conditional inconsistency]\label{thm:splitincons}
The following theory is inconsistent:
\[
 \ZFC+\exists\lambda\bigl(\UEx(\lambda)\wedge\NS(\lambda)\bigr).
\]
In particular, at an ultraexacting $\lambda$ there is no
$\OD_{V_\lambda}$ function $F:Q_\lambda\to\Pow(\lambda)$ such that
$F(q)$ splits $q$ for every $q\in Q_\lambda$.
\end{theorem}

\begin{proof}
Take the critical tail class supplied by Theorem~\ref{thm:nosplit}.
It has no splitter in $\OD(q,p)$ for any $p\in V_\lambda$, contradicting
$\NS(\lambda)$. If $F$ were ordinal definable from a low-rank parameter
$p$, its value at $q$ would be in $\OD(q,p)$ by substitution, giving
the same contradiction.
\end{proof}

This is not a failure of Choice. Since $Q_\lambda$ is a set, Choice
selects a representative for every class. Taking the even entries
of each selected representative gives an unrestricted splitting
function. Thus the inconsistency theorem isolates the definability
restriction, rather than disguising an impossible demand for ordinary
splitting subsets.

\subsection{The new structural package}

Let $\NMPr(\lambda)$ abbreviate the following assertion. There is
$q\in Q_\lambda$ such that, for every $p\in V_\lambda$, the class
$M_{q,p}$ of Definition~\ref{def:hod} contains the normal measure
$U_{q,p}$ of \eqref{eq:U}, and every $a\in q$ generates a
$\PP_{U_{q,p}}^{M_{q,p}}$-generic filter over that class. The resulting
$K_{q,p}$ is independent of $a$ and has the preservation properties
of Theorem~\ref{thm:layer}.

This notation is not a new second-order axiom. Membership in
$M_{q,p}$ is expressed by a fixed first-order formula. Normality and
completeness are relativizations to that formula. Genericity says
that the set $G_a$ meets every dense subset belonging to the class;
that too is a first-order assertion. Name evaluation defines the
corresponding layer. Statements that these classes satisfy individual
axioms of $\ZFC$ follow by the arguments already given, rather than
by quantifying over a truth predicate for a proper class.

\begin{theorem}[Intrinsic structure]\label{thm:package}
\[
 \ZFC\vdash\forall\lambda\,
       \bigl(\UEx(\lambda)\longrightarrow\NMPr(\lambda)\bigr).
\]
With $V_\lambda\subseteq\HOD$ added, one can in addition require that
all the fixed-parameter models and measures coincide and that their
rank segments below $\lambda$ equal those of the ambient universe.
\end{theorem}

\begin{proof}
Apply Theorems~\ref{thm:measure} and \ref{thm:layer}, followed in the
coded case by Corollary~\ref{cor:boundary}.
\end{proof}

Here $\Izero$ means a nontrivial elementary embedding
$i:L(V_{\eta+1})\to L(V_{\eta+1})$ with critical point below
$\eta$. The comparison does not identify $\eta$ with the
$\lambda$ occurring in a different consistency model.

\begin{theorem}[Enhanced $\Izero$ calibration]\label{thm:equicons}
The following theories are equiconsistent:
\begin{align*}
 T_0&=\ZFC+\text{``There is an $\Izero$ embedding''},\\
 T_{\mathrm{layer}}
   &=\ZFC+\exists\lambda\bigl(\UEx(\lambda)
       \wedge V_\lambda\subseteq\HOD\wedge\NMPr(\lambda)\bigr).
\end{align*}
In the second theory one may also include the full low-rank agreement
and common-model assertions of Corollary~\ref{cor:boundary}, and
$\neg\NS(\lambda)$, without changing the conclusion.
\end{theorem}

\begin{proof}
\emph{Forward relative consistency.}
The imported comparison theorem gives
\[
 \Con(T_0)\ \Longrightarrow\
 \Con(\ZFC+\exists\lambda\,\UEx(\lambda)).
\]
In an ultraexacting model use the imported coding forcing of
\cite[Proposition~3.9]{abgl} to retain ultraexactingness of $\lambda$
and obtain $V_\lambda\subseteq\HOD$. This is a relative-consistency
construction, so it gives the corresponding implication between
formal consistency statements, not an assumption that the universe
already has an external generic filter.

In the resulting model, apply the internal $\ZFC$ proof of
Theorem~\ref{thm:package}. It supplies $\NMPr(\lambda)$, the common
model and full rank agreement. Theorem~\ref{thm:splitincons} supplies
$\neg\NS(\lambda)$. Thus the enhanced theory is relatively consistent.

\emph{Reverse relative consistency.}
Forget all the extra assertions of $T_{\mathrm{layer}}$ and retain an
ultraexacting cardinal. The reverse direction of the imported
comparison theorem then gives $\Con(T_0)$.
\end{proof}

\begin{assessment}{What the calibration does and does not establish}
The additional structure proved here costs no consistency strength
\emph{over the already calibrated ultraexacting boundary theory}.
This is not a new proof of the ultraexacting--$\Izero$ comparison and
not a consistency-strength separation. Conversely, the proof supplies
more than an example in a specially chosen Prikry extension:
Theorem~\ref{thm:package} constructs a canonical relative-HOD layer in
\emph{every} universe with an ultraexacting cardinal.
\end{assessment}

\section{The stationary-set gap above the canonical layer}\label{sec:stationary}

This section makes the archive's successor obstruction concrete for
the new models. The general collapse-or-destroy argument was already
present in the synthesis \cite[the successor-stationarity theorem]{synthesis};
the canonical Prikry layer, and its exact role in the application,
are the additional structure here.

Fix $q,p$, and abbreviate
\[
 M=M_{q,p},\qquad K=K_{q,p},\qquad
 \theta=(\lambda^+)^M=(\lambda^+)^K,
 \qquad S=(E_\lambda^\theta)^M.
\]
Thus $S$ is the set of ordinals below $\theta$ whose cofinality is
$\lambda$ in $M$.

\begin{lemma}[The stationary witness]\label{lem:stationary}
$S$ is stationary in $\theta$ in $K$, and
$K\models S\subseteq E_\omega^\theta$.
In every outer model $W\supseteq K$ of $\ZFC$, $S$ fails to reflect at
every $\beta<\theta$ with $\cf^W(\beta)>\omega$.
\end{lemma}

\begin{proof}
In $M$, $\lambda$ is regular and $\theta=\lambda^+$ is regular, so
$E_\lambda^\theta$ is stationary. Explicitly, inside any club in
$\theta$ construct an increasing continuous sequence of length
$\lambda$; its supremum is below $\theta$, is in the club, and has
cofinality $\lambda$. Normal Prikry forcing is $\theta$-cc in $M$,
so it preserves this stationary set. A proof of this chain-condition
preservation fact is included below in Lemma~\ref{lem:cc}.

For each $\alpha\in S$, $M$ has an increasing cofinal map
$h_\alpha:\lambda\to\alpha$. Composing it with a cofinal
$\omega$-sequence in $\lambda$ belonging to $K$ gives a cofinal
$\omega$-sequence in $\alpha$. Thus $\cf^K(\alpha)=\omega$,
and the same is true in every outer $W$.

Now fix $\beta<\theta$ with $\cf^W(\beta)>\omega$. Let
$\mu=\cf^M(\beta)$. Since $\beta<\lambda^+$ in $M$,
$\mu\le\lambda$. If $\mu=\lambda$, composition as above would give
$\cf^W(\beta)=\omega$, impossible. If $\mu\le\omega$, the cofinal
sequence in $M$ would already contradict $\cf^W(\beta)>\omega$.
Consequently $\omega<\mu<\lambda$.

In $M$ choose a strictly increasing continuous cofinal map
$h:\mu\to\beta$. The set
\[
 C=\{h(\xi):\xi<\mu\text{ is a nonzero limit ordinal}\}
\]
is club in $\beta$. At each such $\xi$, continuity gives
$\cf^M(h(\xi))\le\cf^M(\xi)<\lambda$; hence $C\cap S=\varnothing$.
Being closed and unbounded in a fixed ordinal is absolute between
transitive models containing the set, so this same $C$ is a club
witness in $W$. Therefore $S\cap\beta$ is nonstationary in $\beta$.
\end{proof}

For completeness we give the reflection fact used in the application.
Use the usual ultrafilter formulation of strong compactness: for every
cardinal $\tau\ge\delta$ there is a fine $\delta$-complete ultrafilter
on $P_\delta(\tau)$.

\begin{lemma}[Strong compactness reflects countable cofinality]\label{lem:sc}
If $\delta$ is strongly compact and $\theta>\delta$ is regular, every
stationary $T\subseteq E_\omega^\theta$ reflects at some
$\beta<\theta$ with $\omega<\cf(\beta)<\delta$.
\end{lemma}

\begin{proof}
Take the well-founded ultrapower $i:V\to N$ by a fine
$\delta$-complete ultrafilter on $P_\delta(\theta)$. Its seed gives a
set $t\in N$ such that $i``\theta\subseteq t\subseteq i(\theta)$ and
$|t|^N<i(\delta)$. Put $\gamma=\sup i``\theta$.
Since $i(\theta)$ is regular in $N$ and exceeds $i(\delta)$,
$\gamma<i(\theta)$; also $\cf^N(\gamma)<i(\delta)$.
Externally $\cf(\gamma)=\theta$: an increasing cofinal sequence of
length $\theta$ is supplied by $i``\theta$, and a shorter cofinal
set would pull back to an unbounded subset of the regular $\theta$
of smaller size. In particular $\cf^N(\gamma)>\omega$.

We claim $i(T)\cap\gamma$ is stationary in $\gamma$ in $N$.
Let $C\in N$ be any club in $\gamma$. For each $\alpha<\theta$, choose
$g(\alpha)>\alpha$ such that some point of $C$ lies strictly between
$i(\alpha)$ and $i(g(\alpha))$. This is possible because both $C$ and
$i``\theta$ are cofinal in $\gamma$.
Let $D\subseteq\theta$ be the club of limit ordinals closed under $g$.
Pick $\alpha\in D\cap T$. Since $\cf(\alpha)=\omega$ and
$\crit(i)=\delta>\omega$, $i$ is continuous at $\alpha$.
Choose $\alpha_n\nearrow\alpha$; there are points
$d_n\in C$ with
\[
 i(\alpha_n)<d_n<i(g(\alpha_n))<i(\alpha).
\]
Their supremum is $i(\alpha)$, so closure of $C$ gives
$i(\alpha)\in C\cap i(T)\cap\gamma$.
This proves the claim for every club $C$ in $N$.

In $N$, therefore, $\gamma<i(\theta)$ witnesses reflection of $i(T)$
at an ordinal of uncountable cofinality below $i(\delta)$.
Elementarity gives the claimed reflection point below $\theta$ in $V$.
\end{proof}

\begin{theorem}[A canonical collapse-or-destroy alternative]\label{thm:gap}
If the ambient universe $V$ has a strongly compact cardinal
$\delta<\lambda$, then exactly one of the following cardinal cases
holds, with the indicated additional conclusion in the second case:
\begin{enumerate}
\item $|\theta|^V=\lambda$; or
\item $\theta=(\lambda^+)^V$ and $S$ is nonstationary in $V$.
\end{enumerate}
In particular $K\ne V$. There is no representation $V=K[H]$ by a
forcing that is $\theta$-cc in $K$, and no such representation by a
forcing with a dense subset of $K$-cardinality below $\theta$.
\end{theorem}

\begin{proof}
Every ordinal in $[\lambda,\theta)$ has $M$-cardinality $\lambda$,
so has $V$-cardinality $\lambda$ as well. Since $\lambda$ is still a
cardinal in $V$, either $\theta$ also has $V$-cardinality $\lambda$,
or it is precisely $(\lambda^+)^V$. This gives the two cardinal cases.

In the latter case, suppose $S$ remained stationary. The cofinal maps
in Lemma~\ref{lem:stationary} show $S\subseteq E_\omega^\theta$ in $V$.
Lemma~\ref{lem:sc} would make it reflect at an ordinal of uncountable
cofinality. This contradicts the final assertion of
Lemma~\ref{lem:stationary}. Therefore $S$ is nonstationary.

In $K$, $\theta$ is still a cardinal and $S$ is still stationary;
hence $K\ne V$. A $\theta$-cc forcing over $K$ preserves both those
facts, again contradicting the alternative. A dense subset of size
below $\theta$ implies the $\theta$ chain condition, proving the
last assertion.
\end{proof}

This theorem does not refute an ultraexacting cardinal above a
strongly compact cardinal. It locates what the ambient universe must
do \emph{beyond the canonical Prikry layer}. The same argument works
with the relevant stationary-reflection principle in place of strong
compactness. No consistency of the combined large-cardinal hypothesis
is being assumed or established here.

\section{Audit, limitations, and next mathematical questions}\label{sec:audit}

\subsection{Proof audit}

The central implications use two counterexample predicates, not a
scheme of individually chosen embeddings. Uniform correctness is
fixed first; the critical tail class is then extracted from one
witness; only afterward are all low-rank parameters quantified over.
The essential checks are the following.

\begin{center}
\small
\renewcommand{\arraystretch}{1.2}
\begin{tabularx}{\textwidth}{@{}P{.32\textwidth}Y@{}}
\toprule
\textbf{Potential failure}&\textbf{Reason it does not occur here}\\\midrule
A moved ordinal parameter is silently fixed&The minimum parameter rank and the canonical counterexample are uniquely defined from fixed parameters; the original ordinal code is never assumed fixed.\\
The embedding domain is treated as transitive&Only $V_\lambda\subseteq X$ and elementary definable closure are used. The calculation on a countable set uses its enumeration in $X$.\\
A representative is admitted as a parameter&$q$ is a single set parameter, not a list of permitted representatives. The separate bound is $p\in V_\lambda$.\\
The measure is merely external&Equation~\eqref{eq:U} and hereditary definability explicitly prove $U_{q,p}\in M_{q,p}$.\\
Internal completeness is confused with ambient completeness&Only sequences belonging to $M_{q,p}$ are used in the completeness proof. Equation~\eqref{eq:externalincomplete} exhibits external failure.\\
A shifted orbit is assumed normal&Proposition~\ref{prop:shift} gives a concrete counterexample and isolates the special role of $\kappa_0$.\\
Prikry genericity is asserted from cofinality alone&The tail condition is checked for every internal measure-one set, and the dense-set proof of Mathias's criterion is supplied.\\
An inner layer is called a ground of $V$&Only $M_{q,p}\subsetneq K_{q,p}$ is proved to be a forcing-ground relation. The further inclusion into $V$ is not assigned a forcing.\\
The $\Izero$ calibration is presented as new&The comparison and coding inputs are identified separately; only the enlarged canonical structure is proved here.\\
\bottomrule
\end{tabularx}
\end{center}

\subsection{Questions not resolved by these proofs}

\begin{question}
Can ordinary exactingness produce a tail class satisfying the two
tail theorems, or an internal normal measure in some comparably
canonical relative-definability model?
\end{question}
The present proof uses $e\in X$ to internalize the critical sequence
and fix its full tail class. Removing that step would require a new
argument, not merely deletion of ``ultra'' from the hypotheses.

\begin{question}
Under additional hypotheses, can one identify the relation between
$K_{q,p}$ and $V$ more precisely, for example as a specified forcing
extension or by a covering/approximation property?
\end{question}
Theorem~\ref{thm:gap} gives an obstruction to small-chain-condition
representations in the presence of a lower strongly compact cardinal.
It gives no positive quotient forcing.

\begin{question}
How many different normal measures or relative HOD models arise from
critical tail classes of sufficiently correct witnesses?
\end{question}
The archive's many rigid classes do not settle this: recodings can
produce the same model, and even change a normal tail presentation
into a nonnormal one. Counting classes, models, and normal measures
are three distinct problems.

\begin{question}
What additional hypothesis can eliminate the parameter $q$ from the
measure, so that a normal measure on $\lambda$ belongs to ordinary
$\HOD$?
\end{question}
Neither $V_\lambda\subseteq\HOD$ nor the present canonicality argument
eliminates $q$. It is important not to claim ordinary-HOD measurability
from Corollary~\ref{cor:boundary}.

\subsection{Verification and source scope}

The supplied TeX synthesis and the Lean source interfaces were used
as the baseline. No new Lean file was compiled and no machine
verification of the new results is claimed. The ultraexacting witness
interface and the comparison and coding results were checked in the
primary arXiv versions cited below; Mathias's criterion was checked
in Benhamou's treatment. Local Kunen remains the standard imported
inconsistency theorem; its original proof is not reproduced or newly
verified here. The literature check was targeted, not an exhaustive
priority search. In particular, the established critical-sequence/
Prikry connection should not be mistaken for a novelty claim.

The complete paper proof supports the internal normal-measure theorem,
its canonical-layer corollary, and the stated conditional consistency
and inconsistency assertions. Specialist review should focus first on
the uniform reflection convention and on the use of \emph{hereditary}
ordinal definability with a single high-rank set parameter.

\appendix
\section{Prikry genericity and preservation: proofs}\label{sec:prikryproof}

This appendix gives the classical facts used above, with all
constructions performed inside a transitive model $M\models\ZFC$.
Within this appendix temporarily write $\nu$ for its measurable
cardinal and $U\in M$ for its nonprincipal normal $\nu$-complete
ultrafilter. All sequences of sets selected in the lemmas belong to
$M$. A reader can then take $\nu=\lambda$ and $M=M_{q,p}$.
The criterion is due to Mathias \cite{mathias};
\cite[Section~3]{benhamou} is a further reference for these classical
Prikry facts. The arguments below are included to expose the exact
hypotheses needed in the application.

For a finite increasing sequence $s$, write
$A_{>s}=\{\alpha\in A:(\forall\xi\in\ran(s))\ \xi<\alpha\}$.
For the empty sequence, $A_{>s}=A$.

\subsection{Normality gives diagonal thinning and homogeneity}

\begin{lemma}[Finite-stem diagonal intersection]\label{lem:diag}
If $\langle A_s:s\in[\nu]^{<\omega}\rangle\in M$ and every $A_s$ is in
$U$, then
\[
 B=\{\beta<\nu:\forall s\in[\beta]^{<\omega}\ \beta\in A_s\}
\]
belongs to $U$.
\end{lemma}

\begin{proof}
Suppose its complement belongs to $U$. Intersect that complement with
$A_\varnothing$ and remove zero. For each remaining $\beta$, choose a
nonempty finite $s_\beta\subseteq\beta$ such that
$\beta\notin A_{s_\beta}$. Choice in $M$ makes this an internal
function. The function $\beta\mapsto\max(s_\beta)$ is regressive.
Normality gives a set $T\in U$ and $\alpha<\nu$ such that
$\max(s_\beta)=\alpha$ for every $\beta\in T$.

There are fewer than $\nu$ finite subsets of $\alpha+1$.
By $\nu$-completeness,
\[
 T\cap\bigcap\{A_s:s\in[\alpha+1]^{<\omega}\}\in U.
\]
Any $\beta$ in that intersection must belong to $A_{s_\beta}$,
contradicting its selection. Thus $B\in U$.
\end{proof}

\begin{lemma}[Finite-arity homogeneity]\label{lem:homogeneous}
If $r<\nu$ and $F:[\nu]^{<\omega}\to r$ belongs to $M$, there is
$H\in U$ such that $F$ is constant on $[H]^n$ for each fixed
$n<\omega$. The constant may depend on $n$.
\end{lemma}

\begin{proof}
First prove the assertion for a single arity by induction. Arity zero
is trivial. At arity one, exactly one of the fewer than $\nu$ color
classes belongs to $U$: if none did, intersecting their complements
would contradict $\nu$-completeness.

Suppose the result holds for arity $n$ and consider
$F:[\nu]^{n+1}\to r$. For each increasing $n$-tuple $s$, choose a
color $i_s<r$ such that
\[
 A_s=\{\beta>\max(s):F(s\suff\langle\beta\rangle)=i_s\}\in U.
\]
For $n=0$, interpret the bound on $\beta$ as empty. The color exists
by completeness and the ultrafilter property. By the induction
hypothesis, there is $H_0\in U$ on whose $n$-tuples $s$ the value
$i_s$ is constant. Apply Lemma~\ref{lem:diag} to the $A_s$, assigning
$\nu$ to unused arities, and intersect the resulting set with $H_0$.
For any $(n+1)$-tuple from this intersection, its last entry is in
the $A_s$ indexed by the preceding entries. Its color is therefore
the common $i_s$. This completes the induction.

For a coloring of all finite arities, select a homogeneous set for
each $n$ and intersect them. This is an internal countable sequence;
$\omega<\nu$, so $\nu$-completeness preserves membership in $U$.
\end{proof}

\subsection{The strong Prikry lemma}

\begin{lemma}[Strong Prikry lemma]\label{lem:strongprikry}
Let $D\in M$ be dense open in $\PP_U^M$. For each condition $(s,A)$
there are $A^*\subseteq A$ in $U$ and $n<\omega$ such that
\[
 (s\suff t,A^*_{>s\suff t})\in D
 \qquad\text{for every }t\in[A^*]^n.
\]
\end{lemma}

\begin{proof}
For each $t\in[A]^{<\omega}$, color $t$ by one if some condition with
stem $s\suff t$ extending $(s,A)$ belongs to $D$, and by zero
otherwise. When its color is one, choose $B_t\in U$ witnessing
\[
 (s\suff t,B_t)\le(s,A),\qquad (s\suff t,B_t)\in D.
\]
For other indices assign $B_t=\nu$; also assign $\nu$ to finite
sequences outside $A$.

Extend the coloring arbitrarily to finite sequences outside $A$.
Choose $H\subseteq A$ in $U$ homogeneous at each arity for the
coloring, using Lemma~\ref{lem:homogeneous} and then intersecting
with $A$. Let $B\in U$ be the
finite-stem diagonal intersection of all the $B_t$, and put
$A^*=H\cap B$. Whenever $t$ has color one,
$A^*_{>s\suff t}\subseteq B_t$, so openness of $D$ gives
$(s\suff t,A^*_{>s\suff t})\in D$.

By density, some extension of $(s,A^*)$ is in $D$. Its additional
stem is a finite $t\subseteq A^*$, of some length $n$, and its color
is one. Homogeneity at arity $n$ makes every $n$-tuple from $A^*$
have color one. This is the desired conclusion, including the
possible case $n=0$.
\end{proof}

\subsection{Proof of Mathias's criterion}

Let $a=\langle\alpha_n:n<\omega\rangle$ be strictly increasing and
cofinal in $\nu$, possibly external to $M$, and assume it is
eventually in each member of $U$. Form $G_a$ as in
Section~\ref{sec:layers}.

First $G_a$ is a filter. It is closed toward weaker conditions by the
extension definition. Two stems of conditions in $G_a$ are initial
segments of the same sequence, hence comparable. Take the longer
stem and intersect their measure sets, deleting entries below that
stem. This gives a common stronger condition still in $G_a$.

Fix a dense open $D\in M$. For each finite stem $s$, apply
Lemma~\ref{lem:strongprikry} below $(s,\nu_{>s})$, and choose in $M$
a set $A_s\in U$ and a number $n_s$ such that
\[
 (s\suff t,(A_s)_{>s\suff t})\in D
 \quad(t\in[A_s]^{n_s}).
\]
Let $B\in U$ be their finite-stem diagonal intersection.
The sequence $a$ is eventually in $B$. Choose an initial segment
$s$ of $a$ containing all entries before that point. Every subsequent
entry $\beta$ of $a$ is in $B$ and is greater than all entries of $s$,
so the diagonal property gives $\beta\in A_s$.

Let $t$ be the next $n_s$ entries of $a$. The condition
$(s\suff t,(A_s)_{>s\suff t})$ is in $D$. Its stem is an initial
segment of $a$ and all remaining entries are in its measure set,
so it is also in $G_a$. Hence $G_a$ meets every dense open set of $M$,
which proves Lemma~\ref{lem:mathias}. Arbitrary dense sets can be
replaced by their downward closures; a filter meeting that closure
also meets the original dense set by upward closure.

Conversely a Prikry-generic sequence is eventually in every
$A\in U$, since conditions whose measure part is contained in $A$
are dense. Conditions whose stem has an entry above any prescribed
bound are also dense, as are conditions with stem length at least
any prescribed natural number. These verify the familiar necessary
cofinality and tail properties.

\subsection{The ordinary Prikry property and bounded subsets}

\begin{lemma}[Prikry property]\label{lem:prikryproperty}
Every forcing statement is decided by a direct extension of any
given condition.
\end{lemma}

\begin{proof}
Fix a statement $\varphi$ and a condition $(s,A)$.
For each finite $t$ from $A$, assign color zero if some direct
extension with stem $s\suff t$ forces $\varphi$, color one if one
forces $\neg\varphi$, and color two if neither exists.
Colors zero and one cannot both be possible, because conditions with
the same stem have a common direct extension by intersecting their
measure sets.

For each tuple with color zero or one choose a witnessing measure
set $B_t$. Apply finite-arity homogeneity and diagonal thinning,
exactly as in Lemma~\ref{lem:strongprikry}, to obtain $A^*\subseteq A$
in $U$ such that every tuple with color zero or one has its
canonical condition $(s\suff t,A^*_{>s\suff t})$ forcing the assigned
truth value, and the color is constant at each length.

There is a deciding extension of $(s,A^*)$. If its additional stem
has length $n$, the common color at length $n$ is zero or one.
Every condition below $(s,A^*)$ can be further extended to have at
least $n$ additional entries. The resulting condition extends the
canonical condition associated with its first $n$ new entries, so
it forces the same truth value. Thus a dense set below
$(s,A^*)$ forces that value, and $(s,A^*)$ itself forces it.
\end{proof}

\begin{corollary}[No new bounded subsets]\label{cor:nobounded}
A $\PP_U^M$-generic extension has no new subset of any $\eta<\nu$.
\end{corollary}

\begin{proof}
Given a name for a subset of $\eta$ and an initial condition, use the
Prikry property successively for the statements that each
$\xi<\eta$ belongs to the named set. Keep the stem unchanged.
At limit stages intersect the measure sets chosen so far, and at
the end intersect all of them. The recursion has length
$\eta<\nu$, belongs to $M$, and $\nu$-completeness ensures each
intersection remains in $U$. The final direct extension decides
membership for all $\xi<\eta$. The set of the affirmative decisions
is a set of $M$. Such conditions are dense, proving the assertion.
\end{proof}

\subsection{Chain condition, stationary sets, and cardinals}

\begin{lemma}[Chain-condition facts]\label{lem:cc}
In $M$, $\PP_U$ is $\nu^+$-cc. More generally a $\theta$-cc forcing,
where $\theta$ is regular uncountable, preserves stationary subsets
of $\theta$ and all cardinals at least $\theta$.
\end{lemma}

\begin{proof}
For Prikry forcing, there are only $\nu$ finite stems. Any two
conditions with the same stem are compatible, by intersecting their
measure sets. An antichain therefore has size at most $\nu$.

For stationarity, let $\PP$ be $\theta$-cc and let a condition $u$
force that $\dot C$ is club in $\theta$. For each $\alpha<\theta$,
choose a maximal antichain below $u$ deciding a point of $\dot C$
strictly above $\alpha$. There are fewer than $\theta$ possible
chosen points. By regularity their supremum is below $\theta$;
choose a strict upper bound $b(\alpha)<\theta$.
The set of limit ordinals $\gamma<\theta$ closed under $b$ is a
ground-model club $D$. For any $\gamma\in D$ and any $\alpha<\gamma$,
the antichain ensures that $u$ forces a point of $\dot C$ between
$\alpha$ and $\gamma$. Hence $u$ forces $\gamma\in\dot C$ by
closure. A ground-model stationary set meets $D$, so no condition
can force that it misses $\dot C$.

For cardinal preservation, a name for an ordinal value has fewer
than $\theta$ possible values once a maximal deciding antichain is
chosen. Therefore a name for a function with domain a cardinal
$\mu<\xi$, where $\xi\ge\theta$ is a ground-model cardinal, cannot
surject onto $\xi$. If $\xi=\theta$, regularity bounds the union of
$\mu<\theta$ many sets of possible values below size $\theta$.
If $\xi>\theta$, that union has size at most
$\mu\cdot\theta<\xi$. This excludes collapse of any such cardinal.
The same argument bounds the possible range of a short alleged
cofinal sequence in $\theta$, so $\theta$ stays regular as used in
the stationarity assertion.
\end{proof}

\begin{corollary}[All-cardinal preservation]\label{cor:cardinalpres}
Normal Prikry forcing preserves every cardinal of $M$, changes
$\cf(\nu)$ to $\omega$, and preserves $(\nu^+)^M$.
\end{corollary}

\begin{proof}
Cardinals at least $\nu^+$ are preserved by Lemma~\ref{lem:cc}.
Since $\nu$ is strongly inaccessible in $M$, functions between
ordinals below $\nu$ can be coded by bounded subsets of $\nu$.
Corollary~\ref{cor:nobounded} therefore excludes collapse of any
cardinal below $\nu$.

If $\nu$ itself collapsed to some $\mu<\nu$, an injection from $\nu$
into $\mu$ would restrict to an injection from a ground-model
cardinal $\rho$ with $\mu<\rho<\nu$ into $\mu$. That restriction is
coded below $\nu$ and so would already belong to $M$, a
contradiction. Such a $\rho$ exists because $\nu$ is a limit cardinal.
Finally the generic union of stems has order type $\omega$ and is
cofinal in $\nu$, by the dense sets described above.
\end{proof}

All the completeness and homogeneity arguments in this appendix were
internal to $M$. The only external object was the candidate cofinal
sequence in Mathias's criterion, and it was used \emph{after} the
internal dense-set constructions. This is why the appendix applies
even though the tail measure in the ambient universe is explicitly
countably incomplete.

\clearpage
\phantomsection
\addcontentsline{toc}{section}{References}
\begin{thebibliography}{99}
\small
\setlength{\itemsep}{6pt}

\bibitem{synthesis}
\emph{Large cardinals at the inconsistency frontier: a synthesis}.
Second edition, 18 September 2026. Supplied in \nolinkurl{Cardinals3.zip},
\nolinkurl{docs/Large_Cardinals_Synthesis.tex}. In particular the
critical-sequence normalization, minimal-rank-parameter device,
cofinal-orbit rigidity, coded boundary package, and successor-stationarity
obstruction. Archive source, not a publication or independent verification.

\bibitem{plan}
\emph{Large cardinals at the inconsistency frontier: a unified assessment}.
18 September 2026. Supplied in \nolinkurl{Cardinals3.zip},
\nolinkurl{docs/Large_Cardinals_Unified_Report.tex}.

\bibitem{abl}
J.~P.~Aguilera, J.~Bagaria, P.~L\"ucke.
\emph{Large cardinals, structural reflection, and the HOD Conjecture}.
\arxiv{2411.11568v4}, revised 16 September 2025.
Definition~3.3 and Lemma~3.2 for the ultraexacting witness interface.
Primary version checked 18 September 2026.

\bibitem{abgl}
J.~P.~Aguilera, J.~Bagaria, G.~Goldberg, P.~L\"ucke.
\emph{Large cardinals beyond HOD}.
\arxiv{2509.10254v1}, 12 September 2025.
Theorem~A (equiconsistency), Proposition~3.9 (coding), and
Section~5.2 (comparison with an iteration-based Prikry construction).
Primary version checked 18 September 2026.

\bibitem{kunen}
K.~Kunen.
Elementary embeddings and infinitary combinatorics.
\emph{Journal of Symbolic Logic} \textbf{36}(3) (1971), 407--413.
\doi{10.2307/2269948}.
The local Kunen theorem is the standard imported inconsistency input,
as in the supplied synthesis.

\bibitem{mathias}
A.~R.~D.~Mathias.
On sequences generic in the sense of Prikry.
\emph{Journal of the Australian Mathematical Society}
\textbf{15}(4) (1973), 409--414.
\doi{10.1017/S1446788700028755}.
Attribution of the genericity criterion; a proof is included in this report.

\bibitem{benhamou}
T.~Benhamou.
\emph{Prikry forcing and tree Prikry forcing of various filters}.
\arxiv{1801.04424v2}, revised 25 September 2018.
Journal version: \doi{10.1007/s00153-019-00660-3}.
Section~3, especially Lemma~3.16 and Theorem~3.17.
Version dates follow the arXiv submission record, not the date inserted
by an experimental HTML rendering.

\end{thebibliography}
\end{document}
